\documentclass[11pt,a4paper]{article}

\usepackage[T1]{fontenc}
\usepackage[utf8]{inputenc}
\usepackage[italian]{babel}
\usepackage[margin=25mm]{geometry}
\usepackage{amsmath,amssymb}
\usepackage{booktabs}
\usepackage{longtable}
\usepackage{array}
\usepackage{siunitx}
\usepackage{enumitem}
\usepackage{fancyhdr}
\usepackage[hidelinks]{hyperref}

\sisetup{output-decimal-marker={,}, per-mode=symbol, detect-all}

\pagestyle{fancy}
\fancyhf{}
\lhead{\small EN 50160 --- Nota riassuntiva}
\rhead{\small\thepage}
\renewcommand{\headrulewidth}{0.4pt}

\newcommand{\Uc}{U_{\mathrm{c}}}
\newcommand{\Un}{U_{\mathrm{n}}}

\title{\textbf{EN 50160 --- Caratteristiche della tensione fornita\\dalle reti pubbliche di distribuzione}\\[2mm]
\large Nota tecnica riassuntiva}
\author{}
\date{\today}

\begin{document}
\maketitle
\thispagestyle{fancy}

\begin{center}
\fbox{\begin{minipage}{0.92\textwidth}
\small\textit{Documento di sintesi a uso interno. I valori riportati sono ripresi dalla
letteratura tecnica corrente e hanno carattere indicativo: per ogni verifica formale
di conformità fa fede esclusivamente il testo ufficiale della norma
EN 50160:2022 (+A1:2025) e i relativi recepimenti nazionali (CEI EN 50160 in Italia).}
\end{minipage}}
\end{center}

\vspace{4mm}

\section{Scopo e campo di applicazione}

La EN 50160 descrive le \emph{caratteristiche attese} della tensione al punto di
consegna dell'utente nelle reti pubbliche di distribuzione in corrente alternata,
in condizioni normali di esercizio. È una norma di \emph{prodotto-elettricità}:
definisce che cosa il cliente può aspettarsi dalla rete, non che cosa il cliente
può iniettare in rete.

\begin{itemize}[itemsep=2pt]
  \item \textbf{Edizione corrente:} EN 50160:2022, approvata da CENELEC nel novembre 2022,
        che sostituisce la EN 50160:2010 e tutte le sue varianti
        (A1:2015, A2:2019, A3:2019); successivamente integrata dall'emendamento A1:2025.
  \item \textbf{Livelli di tensione coperti:} BT ($\Un \le \SI{1}{kV}$),
        MT ($\SI{1}{kV} < \Un \le \SI{36}{kV}$), AT e --- novità dell'edizione 2022 ---
        AAT. L'edizione 2022 dedica un capitolo separato a ciascun livello.
  \item \textbf{Esclusioni:} reti industriali e private, impianti in condizioni eccezionali
        (guasti di forza maggiore, condizioni meteorologiche estreme, azioni di terzi,
        distacchi per emergenza di sistema, atti di autorità), reti in isola non pubbliche.
  \item \textbf{Natura statistica:} i limiti non sono soglie istantanee, ma percentili
        valutati su intervalli di osservazione tipicamente settimanali
        (95\,\% o 99\,\% dei valori aggregati). Ciò è il punto in cui la norma si
        differenzia sostanzialmente dai livelli di compatibilità EMC.
\end{itemize}

\subsection*{Definizioni essenziali}

\begin{description}[style=nextline,leftmargin=0pt,itemsep=2pt]
  \item[Tensione nominale $\Un$] valore con cui il sistema è designato
        (es. \SI{230}{V} fase--neutro, \SI{400}{V} concatenata in BT).
  \item[Tensione dichiarata $\Uc$] normalmente coincide con $\Un$; può differire se
        concordata contrattualmente tra distributore e utente. In MT/AT è il
        riferimento usato per i limiti.
  \item[Condizioni normali di esercizio] esercizio in cui il carico è alimentato,
        con manovre di rete ed eliminazione automatica dei guasti, in assenza di
        eventi eccezionali.
  \item[Valore aggregato] valore efficace su \SI{10}{\minute} (armoniche, squilibrio,
        ampiezza) o su 10/12 cicli, secondo le regole di aggregazione della
        IEC/EN 61000-4-30.
\end{description}

\section{Frequenza di alimentazione}

Valore misurato come media su \SI{10}{\second}, riferito a $f_\mathrm{n}=\SI{50}{Hz}$.

\begin{table}[h]
\centering
\begin{tabular}{lll}
\toprule
\textbf{Tipo di rete} & \textbf{Intervallo} & \textbf{Percentile / periodo} \\
\midrule
Interconnessa & \SI{50}{Hz}\,$\pm\,1\%$ (\SIrange{49.5}{50.5}{Hz}) & 99,5\,\% dell'anno \\
Interconnessa & $+4\%\,/\,-6\%$ (\SIrange{47}{52}{Hz}) & 100\,\% del tempo \\
Non interconnessa (isola) & \SI{50}{Hz}\,$\pm\,2\%$ & 95\,\% della settimana \\
Non interconnessa (isola) & \SI{50}{Hz}\,$\pm\,15\%$ & 100\,\% del tempo \\
\bottomrule
\end{tabular}
\end{table}

\noindent
La finestra \SIrange{47}{52}{Hz} è il riferimento pratico per il campo di aggancio
di un PLL di rete e per la logica di protezione di interfaccia (che tuttavia segue
le norme di connessione, non la EN 50160).

\section{Ampiezza della tensione di alimentazione}

Valutata su valori efficaci aggregati su \SI{10}{\minute}, su un periodo di osservazione
di una settimana.

\begin{table}[h]
\centering
\begin{tabular}{lll}
\toprule
\textbf{Livello} & \textbf{Limite} & \textbf{Criterio} \\
\midrule
BT / MT & $\Uc \pm 10\%$ & 95\,\% dei valori \SI{10}{\minute} settimanali \\
BT & $+10\%\,/\,-15\%$ & 100\,\% dei valori \SI{10}{\minute} settimanali \\
AT / AAT & valori indicativi o contrattuali & secondo accordo con il gestore \\
\bottomrule
\end{tabular}
\end{table}

\noindent
Per le reti non interconnesse e per alcune situazioni particolari sono ammesse
deroghe nazionali (Allegato E, ``A-deviations'').

\section{Variazioni rapide di tensione e flicker}

\begin{itemize}[itemsep=2pt]
  \item \textbf{Variazioni rapide (RVC):} in BT e MT generalmente inferiori al 5\,\% di $\Un$;
        in condizioni particolari sono ammessi eventi occasionali fino al 10\,\% di $\Un$.
        L'edizione 2022 dettaglia i valori indicativi nell'Allegato B.
  \item \textbf{Severità del flicker a lungo termine:} $P_{lt} \le 1$ per il 95\,\% del tempo
        su una settimana, in BT e MT ($P_{lt}$ calcolato su blocchi di \SI{2}{\hour}
        a partire dai $P_{st}$ su \SI{10}{\minute}, secondo IEC 61000-4-15).
\end{itemize}

\section{Buchi di tensione e interruzioni}

Fenomeni imprevedibili e in larga misura casuali: la norma non fissa limiti prescrittivi
ma fornisce \emph{valori indicativi} basati su statistiche di rete, con classificazione
per profondità e durata (tabella tipo UNIPEDE).

\begin{table}[h]
\centering
\small
\begin{tabular}{p{0.32\textwidth}p{0.60\textwidth}}
\toprule
\textbf{Fenomeno} & \textbf{Valori indicativi} \\
\midrule
Buchi di tensione ($U < 90\%\,\Uc$, durata \SIrange{10}{60000}{ms}) &
da alcune decine fino a circa un migliaio di eventi all'anno; la maggior parte
con durata $<\SI{1}{\second}$ e profondità $<60\%$. \\
\addlinespace
Interruzioni brevi ($\le \SI{3}{\minute}$) &
tipicamente da alcune decine a alcune centinaia all'anno; circa il 70\,\% con
durata inferiore a \SI{1}{\second}. \\
\addlinespace
Interruzioni lunghe ($> \SI{3}{\minute}$) &
tipicamente meno di 10 fino a 50 eventi all'anno, con forte dipendenza dalla
struttura di rete (aerea o in cavo) e dall'area geografica. \\
\bottomrule
\end{tabular}
\end{table}

\noindent
Le interruzioni programmate, preavvisate all'utente, sono contabilizzate separatamente.

\section{Sovratensioni}

\begin{itemize}[itemsep=2pt]
  \item \textbf{Temporanee a frequenza industriale (TOV):} originate tipicamente da
        guasti a monte; in BT generalmente non superano \SI{1.5}{kV} efficaci, con
        entità dipendente dallo stato del neutro della rete MT.
  \item \textbf{Transitorie:} impulsive o oscillatorie, da manovra o da fulminazione;
        in BT generalmente inferiori a \SI{6}{kV} di picco, con tempi di salita
        da microsecondi a millisecondi. Il coordinamento dell'isolamento e la scelta
        degli SPD si basano sulle categorie di sovratensione della EN 60664-1
        e sulla EN 62305, non sulla EN 50160.
\end{itemize}

\section{Squilibrio della tensione}

Rapporto tra componente di sequenza inversa e componente di sequenza diretta,
su valori aggregati a \SI{10}{\minute}:

\begin{center}
\begin{tabular}{ll}
\toprule
\textbf{Limite} & \textbf{Criterio} \\
\midrule
$u_2 \le 2\%$ & 95\,\% dei valori \SI{10}{\minute} su una settimana \\
$u_2 \le 3\%$ & ammesso in alcune aree con forte presenza di carichi monofase \\
\bottomrule
\end{tabular}
\end{center}

\noindent
Nelle reti con neutro distribuito può essere presente anche componente omopolare,
non oggetto di limite esplicito nella norma.

\section{Armoniche e interarmoniche}

Valori efficaci delle singole armoniche riferiti alla fondamentale, aggregati su
\SI{10}{\minute}: il limite deve essere rispettato dal 95\,\% dei valori settimanali.
La distorsione totale, calcolata fino al 40\textsuperscript{o} ordine, deve rispettare
$\mathrm{THD}_U \le 8\%$ in BT e MT.

\begin{table}[h]
\centering
\small
\begin{tabular}{cc cc cc}
\toprule
\multicolumn{2}{c}{\textbf{Dispari non multiple di 3}} &
\multicolumn{2}{c}{\textbf{Dispari multiple di 3}} &
\multicolumn{2}{c}{\textbf{Pari}} \\
\cmidrule(lr){1-2}\cmidrule(lr){3-4}\cmidrule(lr){5-6}
$h$ & $U_h/U_1$ & $h$ & $U_h/U_1$ & $h$ & $U_h/U_1$ \\
\midrule
5  & 6,0\,\%  & 3  & 5,0\,\%  & 2         & 2,0\,\% \\
7  & 5,0\,\%  & 9  & 1,5\,\%  & 4         & 1,0\,\% \\
11 & 3,5\,\%  & 15 & 0,5\,\%  & 6\dots24  & 0,5\,\% \\
13 & 3,0\,\%  & 21 & 0,5\,\%  &           &         \\
17 & 2,0\,\%  &    &          &           &         \\
19 & 1,5\,\%  &    &          &           &         \\
23 & 1,5\,\%  &    &          &           &         \\
25 & 1,5\,\%  &    &          &           &         \\
\bottomrule
\end{tabular}
\caption*{\small Valori indicativi delle singole armoniche di tensione in BT e MT.
Per gli ordini superiori al 25\textsuperscript{o} la norma storicamente non fissava
valori, essendo generalmente ridotti e poco prevedibili per effetto delle risonanze
di rete; l'edizione 2022 amplia la trattazione degli ordini elevati.}
\end{table}

\begin{itemize}[itemsep=2pt]
  \item \textbf{Interarmoniche:} livelli in crescita per effetto di convertitori e
        di generazione distribuita; storicamente trattate in modo informativo,
        con attenzione all'interazione con i sistemi di telecomando centralizzato.
  \item \textbf{Banda \SIrange{2}{150}{kHz} (``supraarmoniche''):} introdotta con
        l'emendamento A2:2019 e recepita nell'edizione 2022 tra gli
        ``altri fenomeni'' (Allegato C). È il punto di contatto diretto con
        l'emissione condotta dei convertitori a commutazione ad alta frequenza,
        in particolare con dispositivi SiC/GaN.
\end{itemize}

\section{Tensioni di segnale sulla rete (mains signalling)}

Segnali di telecomando (ripple control, carrier, marcatori di rete) sovrapposti
alla tensione di alimentazione. Il criterio è espresso come curva di ampiezza
ammessa in funzione della frequenza, verificata sul 99\,\% di un giorno con
valori medi su \SI{3}{\second}: ordine di grandezza pari a circa il 9\,\% di $\Uc$
alle frequenze più basse (centinaia di Hz), con andamento decrescente fino a
qualche punto percentuale nella banda dei kHz e a valori dell'ordine dell'1\,\%
oltre i \SI{10}{kHz}.

\section{Misura e valutazione della conformità}

\begin{itemize}[itemsep=2pt]
  \item Metodi di misura secondo \textbf{IEC/EN 61000-4-30}, classe A per la
        valutazione contrattuale e normativa, classe S per il monitoraggio statistico.
  \item Intervallo base 10 cicli (\SI{50}{Hz}), aggregazione a \SI{150}{\second},
        \SI{10}{\minute} e \SI{2}{\hour}.
  \item Periodo di osservazione: una settimana per la maggior parte dei parametri,
        un anno per la frequenza.
  \item Marcatura dei dati (\emph{flagging}): i valori acquisiti durante buchi,
        sovratensioni o interruzioni sono marcati ed esclusi dalle statistiche
        degli altri parametri, per evitare doppi conteggi.
\end{itemize}

\section{Relazione con le altre norme}

\begin{table}[h]
\centering
\small
\begin{tabular}{p{0.30\textwidth}p{0.62\textwidth}}
\toprule
\textbf{Norma} & \textbf{Ruolo} \\
\midrule
EN 50160 & Caratteristiche della tensione fornita (\emph{che cosa arriva} all'utente). \\
EN 61000-2-2 / -2-4 / -2-12 & Livelli di compatibilità in BT pubblica, in impianti
industriali e in MT. Sono i valori di riferimento per la progettazione EMC. \\
EN 61000-3-2 / -3-12 & Limiti di emissione di corrente armonica per apparecchi
$\le \SI{16}{A}$ e \SIrange{16}{75}{A} per fase. \\
IEC/TS 61000-3-4, IEC/TR 61000-3-6 & Valutazione dell'emissione di apparecchi e
impianti di grande potenza, allocazione delle quote di emissione al PCC. \\
EN 61000-4-30 & Metodi di misura della qualità della tensione. \\
CEI 0-16 / CEI 0-21 (Italia) & Regole di connessione alle reti MT/AT e BT: definiscono
gli obblighi lato utente (protezione di interfaccia, campi di frequenza e tensione,
LVFRT), complementari alla EN 50160. \\
\bottomrule
\end{tabular}
\end{table}

\section{Implicazioni per la progettazione di convertitori connessi alla rete}

\begin{enumerate}[itemsep=3pt]
  \item \textbf{Dimensionamento del campo di tensione.} La banda
        $\Uc\,+10\%/-15\%$ (100\,\% del tempo in BT), sommata alla caduta sul filtro
        e all'ondulazione, determina la tensione minima di bus DC per mantenere il
        controllo di corrente in condizioni di massima tensione di rete e, in
        contrapposizione, la tensione massima di blocco dei dispositivi.
  \item \textbf{PLL e stima di frequenza.} Il campo \SIrange{47}{52}{Hz} definisce il
        range di aggancio; la banda passante del PLL va scelta come compromesso tra
        reiezione dello squilibrio (componente $2\omega$ dovuta a $u_2$ fino al 2--3\,\%)
        e prontezza durante i buchi. Uso pratico di DSOGI-PLL o DDSRF-PLL.
  \item \textbf{Squilibrio.} Con $u_2$ fino al 2--3\,\%, un controllo di sola sequenza
        diretta genera ripple di potenza a \SI{100}{Hz} sul bus DC; se il ripple non
        è accettabile, occorre un anello dedicato di sequenza inversa o un controllo
        risonante multi-frequenza.
  \item \textbf{Armoniche di tensione preesistenti.} Un $\mathrm{THD}_U$ fino all'8\,\%
        con quinta al 6\,\% è un disturbo di ingresso, non un errore di controllo:
        richiede risonanti a $6k\pm1$ nel riferimento sincrono oppure una struttura
        di reiezione a banda larga, e va considerato nella verifica di stabilità
        del filtro LCL con impedenza di rete variabile.
  \item \textbf{Buchi e interruzioni brevi.} La statistica dei buchi
        (in maggioranza $<\SI{1}{\second}$ e $<60\%$ di profondità) definisce il
        requisito realistico di ride-through, l'energia da immagazzinare nel bus DC
        e la logica di ripresa sincronizzata.
  \item \textbf{Banda \SIrange{2}{150}{kHz}.} Con dispositivi a banda larga e frequenze
        di commutazione elevate, l'emissione in questa banda va valutata già in fase
        di progetto del filtro EMI: è l'area in cui la normativa è in evoluzione e in
        cui l'interazione tra più convertitori sullo stesso feeder è meno prevedibile.
  \item \textbf{Distinzione fondamentale.} La EN 50160 non è un limite di emissione:
        la conformità dell'apparecchio si dimostra rispetto alle EN 61000-3-x e alle
        regole di connessione nazionali. La EN 50160 fornisce invece le
        \emph{condizioni ambientali} da assumere in ingresso al progetto.
\end{enumerate}

\section*{Riferimenti}

\begin{itemize}[itemsep=2pt]
  \item EN 50160:2022, \textit{Voltage characteristics of electricity supplied by
        public electricity networks}, CENELEC, con emendamento A1:2025.
  \item IEC/EN 61000-4-30, \textit{Testing and measurement techniques ---
        Power quality measurement methods}.
  \item IEC/EN 61000-4-15, \textit{Flickermeter --- Functional and design specifications}.
  \item EN 61000-2-2, EN 61000-2-4, EN 61000-2-12, \textit{Compatibility levels}.
  \item CEI 0-16 e CEI 0-21, regole tecniche di connessione alle reti di distribuzione
        italiane.
\end{itemize}

\end{document}
